%%%%%%%%%%%%%%%%%%%%%%%%%%%%%%%%%%%%%%%%%%%%%%%%%%%%%%%%%%%%%%
%                    CHAPTER 1: INTRODUCTION
%%%%%%%%%%%%%%%%%%%%%%%%%%%%%%%%%%%%%%%%%%%%%%%%%%%%%%%%%%%%%%

\chapter{Introduction}
\label{chap:intro}

\section{Project Overview}
In today's highly competitive software development industry, technical recruitment processes evaluate candidates across diverse domains including core language mechanics, modern frontend libraries, backend distributed systems, relational databases, and enterprise architectural design. Engineering students and job seekers routinely face multi-stage technical interviews requiring rigorous theoretical understanding, rapid problem-solving skills, and the ability to articulate structured, in-depth architectural solutions.

\textbf{OmniFAQ AI} is a comprehensive, modern web-based technical interview training platform engineered to streamline and elevate candidate preparation. Built upon a decoupled client-server architecture, the platform bridges the gap between static study material and dynamic, real-world candidate needs by pairing a rich, pre-loaded repository of domain-specific interview questions with the real-time synthesis power of \textbf{Google Gemini Generative AI}. The application enables learners to search, filter by adaptive difficulty levels (Easy, Medium, Hard), practice conceptual answers with automated keyword feedback, and dynamically expand the platform's repository on-demand.

\section{Motivation and Background}
With the rapid evolution of modern web frameworks, cloud technologies, and developer tooling, software engineering curricula frequently lag behind current industry interview expectations. Candidates typically spend hundreds of hours navigating dozens of disconnected resources, including competitive coding repositories, community forums, medium articles, and fragmented documentation. 

Furthermore, existing question-and-answer repositories are inherently static: if a student is interviewing for a niche framework (e.g., GraphQL, Docker, Kubernetes, WebSockets) that has not been explicitly cataloged by the site administrator, the user receives zero guidance. Recent advancements in Large Language Models (LLMs), particularly Google's Gemini family of generative models, present a transformative opportunity to convert static educational portals into self-learning, adaptive environments capable of generating domain-accurate interview questions and gold-standard model answers on the fly.

\section{Problem Statement}
Contemporary technical interview preparation suffers from four foundational bottlenecks:
\begin{enumerate}
    \item \textbf{Fragmentation of Learning Resources:} Candidates are forced to bounce between competitive programming websites, blogs, and community forums without a unified repository covering theory, architectural questions, and code semantics.
    \item \textbf{Lack of Comprehensive Model Explanations:} Many portals supply interview questions without high-scoring, structured technical answers that demonstrate depth of understanding.
    \item \textbf{Static, Non-Adaptive Repositories:} Conventional FAQ systems cannot adapt dynamically when learners query emerging, modern, or specialized developer topics.
    \item \textbf{Zero Telemetry on Student Knowledge Gaps:} Educators and platform administrators lack empirical data regarding which topics students search for most frequently and where learners encounter unresolved queries.
\end{enumerate}

\section{Objectives of the System}
The primary objectives of the OmniFAQ AI project are as follows:
\begin{itemize}
    \item \textbf{Unified Multi-Domain Question Bank:} To engineer a curated repository covering five core technical disciplines: JavaScript, React, Node.js, SQL, and System Design.
    \item \textbf{On-Demand Generative AI Expansion:} To integrate the Google Gemini API to dynamically generate 3--4 rigorous interview questions with model answers whenever an unresolved topic is searched.
    \item \textbf{Self-Learning Database Persistence:} To implement an automated ingestion pipeline where AI-generated questions are sanitized and permanently written into the database, benefiting all future platform visitors without recurring API costs.
    \item \textbf{Debounced Real-Time Search Engine:} To create a fast, client-side debounced search interface with weighted relevance scoring across question titles, keyword tags, and model answers.
    \item \textbf{Interactive Candidate Practice Engine:} To supply an active recall practice module that analyzes a user's typed response against expected technical keywords and provides real-time coverage scores.
    \item \textbf{Administrative Analytics and Gap Detection:} To build an administrative telemetry dashboard that logs all user searches, computes resolution rates, and surfaces high-priority unresolved queries for curriculum refinement.
    \item \textbf{Robust Dual-Driver Database Deployment:} To design an environment-agnostic backend supporting cloud MySQL as well as automatic, zero-configuration failover to an embedded SQLite engine.
\end{itemize}

\section{Project Scope}
The scope of OmniFAQ AI encompasses full-stack software development, REST API design, cryptographic security, database normalization, prompt engineering, containerization, and cloud deployment:
\begin{itemize}
    \item \textbf{Target Audience:} Computer Science undergraduate students, bootcamp graduates, junior-to-mid level software engineers, technical mentors, and academic institutions.
    \item \textbf{Deployment Scope:} Containerized using Docker for zero-dependency deployment across cloud PaaS providers (e.g., Render, Railway) as well as traditional LAMP/XAMPP local environments.
    \item \textbf{Security Boundaries:} Complete protection against SQL Injection via PDO parameterized queries, secure user authentication with Bcrypt password hashing, and role-based access control.
\end{itemize}

\section{Report Organization}
The remainder of this report is organized as follows:
\begin{itemize}
    \item \textbf{Chapter 2 (Literature Survey and Technology Stack):} Reviews existing interview portals and details the core technologies comprising the frontend, backend, database, AI engine, and deployment infrastructure.
    \item \textbf{Chapter 3 (System Design and Architecture):} Elucidates the three-tier system architecture, Entity-Relationship (ER) model, Data Flow Diagrams (DFDs), and RESTful API endpoint specifications.
    \item \textbf{Chapter 4 (Implementation and Key Features):} Provides in-depth implementation mechanics of debounced search, generative AI integration, practice answer evaluation, and admin telemetry.
    \item \textbf{Chapter 5 (Testing, Deployment, and Results):} Documents test cases, dual-database failover validation, Docker containerization, cloud deployment on Render, and concludes with future directions.
\end{itemize}
